% Technical Abstract. <=2 pages, self-contained.
% CUED REQUIREMENT: must include name, College, project title -- this exact
% copy is also archived separately by the Department.
% Write LAST (5/30) once all numbers are settled.

\section*{Technical Abstract}

\noindent Project title: Directional Reflectance Profile Analysis for
Laid-Line Detection in Historical Paper\\
Author: Zhanfeng Zhou\\
College: TODO\\
Supervisor: TODO\\[1em]

% --- Body of the abstract (write last, target ~2 pages) ------------------
% Paragraph 1: problem context. Handmade paper carries a mould fingerprint;
% laid-line density and chain-line spacing are forensic identifiers used
% by codicologists; current practice is manual; MSI/DRP image stacks now
% exist at scale but no robust automated pipeline turns them into the
% required quantitative data.
%
% Paragraph 2: approach. A five-stage DRP-based pipeline that estimates
% laid-line orientation automatically from the per-pixel directional
% reflectance, then refines that estimate patchwise, then runs a
% spectral-power detector against the enhanced grayscale image, and
% finally reports a small bundle of evaluation metrics (R^2 of harmonic
% fit, self-contrast z-score, split-half stability, wire FWHM
% distribution).
%
% Paragraph 3: phantom validation. A synthetic line generator with
% controllable period, SNR, orientation and wire-width gives the
% benchmark for detector accuracy under known ground truth; period
% recovery within plus/minus 2.5\% across 10--50 px; essentially
% SNR-invariant down to SNR = 2.
%
% Paragraph 4: real-folio results. Pipeline applied to a benchmark of
% nine folios drawn from four manuscript stocks in the Cambridge UL
% collection, with two folios additionally counted by hand to provide
% calibrated ground truth. Sub-1\% density error on the two calibrated
% folios; the remaining seven show a systematic deviation from the
% library spreadsheet records that is attributed to stock-versus-folio
% mismatch rather than to detector failure.
%
% Paragraph 5: conclusions. The pipeline reaches an accuracy that is in
% principle sufficient to discriminate mould pairs. Failure modes are
% characterised (octave aliasing at high density; polarity-flip for
% transmitted-light samples). Future work covers automatic polarity
% detection, joint chain-line analysis, and mould-mate matching.

\newpage
